ここでは高被引用テーマの変化を、高被引用でない論文の研究テーマ（または引用テーマ）と比較する。
各累積世代、各グループごとの最頻出研究テーマと引用テーマを表\ref{tab:topmesh}に示す。これらのうちG1の研究テーマ（高被引用テーマ）にのみにバラエティが確認される。
次に、各累積世代の研究テーマと引用テーマの多様性を見るために各グループのエントロピーの変化を観察した。研究テーマ、引用テーマ、研究テーマと引用テーマ対の出現数のエントロピーを、それぞれ図\ref{fig:arr2}に示す。計数の方法は研究テーマと引用テーマはそれぞれ論文ごとの分数カウントを採用する。研究テーマ対引用テーマのペアについては、ひとつの参照ごとに研究テーマと引用テーマの2部グラフを構成し、エッジの重みの合計が1となる分数カウントを採用する。図17からは高被引用論文グループのみがエントロピーが低く、その他のグループ間ではそれほど差がないことがわかる。なお、1946-1950年の累積世代に関してはMeSHディスクリプターの付与数が十分でなかったため除外している。
また、研究者の引用行動は必ずしも当該論文と同じテーマを引用するものではない。この転換の程度を見るために、グループごとに研究テーマと引用テーマのコサイン距離を観察した。この結果を図18に示す。こちらも高被引用グループのみが他とは異なる様相を示した。また、高被引用グループは1946-2010年の累積世代以降、転換の程度は減少傾向にあり、他のグループとの差はなくなりつつあるように推察される。
さらに、各累積世代間に対して研究テーマの転換の程度を確認するために、各グループの累積世代間の研究テーマのコサイン距離を観察した。累積世代間では論文集合のメンバーが異なるので論文集合全体のテーマの重み付カウントを用いた。この結果を図19に示す。これによれば累積世代間のテーマの転換の程度は一定の傾向を示さなかったが、グループ2のみが常に低い転換の程度を保っている。
以上より、高被引用論文には次の特徴があると考えられる：(1)高被引用論文は異なるテーマから引用されやすい。このことは表6の最頻出テーマから予想され、図18に示す転換の程度もこれを裏付けている。(2)高被引用論文は研究テーマが集中的であり、それを引用する論文の研究テーマもまた集中的である。このことは図17から示唆される。(3)高被引用論文のテーマの多様性は、共時的には低いが通時的には必ずしも低くないと思われる。このことは、図17では他のグループに比べ低いエントロピーが観察されるにもかかわらず、図19では他のグループと変わらない程度の転換が観察されることから推察される。
